\documentclass{article}
\usepackage{amsmath,amssymb,amsthm}
\usepackage{algorithm}
\usepackage{algpseudocode}
\usepackage{enumitem}
\newtheorem{theorem}{Theorem}
\newtheorem{lemma}{Lemma}
\newtheorem{assumption}{Assumption}
\newtheorem{definition}{Definition}
\newtheorem{corollary}{Corollary}
\newcommand{\E}{\mathbb{E}}
\newcommand{\R}{\mathbb{R}}

\begin{document}

\section*{Algorithm Name}
\textbf{SAGA} (Stochastic Average Gradient Augmented) for finite‑sum minimization.

\section*{Mathematical Formulation}
We consider the finite‑sum problem  

\[
\min_{x\in \R^{d}} f(x)\quad\text{with}\quad 
f(x):=\frac{1}{n}\sum_{i=1}^{n}f_i(x),
\]

where each component function \(f_i:\R^{d}\to\R\) is convex and \(L\)-smooth.
For each \(i\) we keep a \emph{memory point} \(\phi_i\in\R^{d}\) and its gradient
\(g_i:=\nabla f_i(\phi_i)\).  
Define the stored average gradient  

\[
\bar g :=\frac{1}{n}\sum_{i=1}^{n} g_i .
\]

At iteration \(k\) a single index \(i_k\) is drawn uniformly from \(\{1,\dots ,n\}\).
The update consists of three parts:

\[
\boxed{
\begin{aligned}
x_{k+1}&=x_k-\alpha\Bigl(\nabla f_{i_k}(x_k)-\nabla f_{i_k}(\phi_{i_k})+\bar g\Bigr),\\[2mm]
\phi_{i_k}&\gets x_k,\\[1mm]
g_{i_k}&\gets \nabla f_{i_k}(x_k),\\[1mm]
\bar g &\gets \bar g+\frac{1}{n}\bigl(g_{i_k}-\nabla f_{i_k}(\phi_{i_k}^{\text{old}})\bigr).
\end{aligned}}
\]

All other \(\phi_i\) and \(g_i\) remain unchanged.  
The step size \(\alpha>0\) is a constant (or a diminishing schedule) chosen
according to the smoothness constant \(L\).

\section*{Pseudocode}
\begin{algorithm}[H]
\caption{SAGA}
\begin{algorithmic}[1]
\Require Number of component functions \(n\), step size \(\alpha\), max‑iterations \(T\)
\State Initialise \(x_0\in\R^{d}\) arbitrarily
\For{$i=1,\dots ,n$}
    \State \(\phi_i^{0}\gets x_0\)
    \State \(g_i^{0}\gets \nabla f_i(\phi_i^{0})\)
\EndFor
\State \(\bar g^{0}\gets \frac{1}{n}\sum_{i=1}^{n} g_i^{0}\)
\For{$k=0$ \textbf{to} $T-1$}
    \State Sample \(i_k\sim\text{Uniform}\{1,\dots ,n\}\)
    \State Compute stochastic gradient \(v_k:=\nabla f_{i_k}(x_k)-g_{i_k}^{k}+\bar g^{k}\)
    \State Update iterate \(x_{k+1}\gets x_k-\alpha v_k\)
    \State \textbf{Update memory:}
    \State \( \bar g^{k+1}\gets \bar g^{k}+ \frac{1}{n}\bigl(\nabla f_{i_k}(x_k)-g_{i_k}^{k}\bigr)\)
    \State \(g_{i_k}^{k+1}\gets \nabla f_{i_k}(x_k),\qquad 
           \phi_{i_k}^{k+1}\gets x_k\)
    \State For all \(j\neq i_k\): \(g_{j}^{k+1}=g_{j}^{k},\;\phi_{j}^{k+1}=\phi_{j}^{k}\)
\EndFor
\State \Return \(x_T\)
\end{algorithmic}
\end{algorithm}

\section*{Dry Run Example}
Minimize the scalar quadratic  

\[
f(w)=\frac{1}{1}\bigl(w-3\bigr)^{2},\qquad \nabla f(w)=2(w-3).
\]

Here \(n=1\); the table contains a single entry \(\phi_1\) and its gradient.
Set \(w_0=0\), step size \(\alpha=0.1\).

\begin{center}
\begin{tabular}{c|c|c|c}
Iteration $k$ & $\phi_1^{k}$ & $v_k$ (update direction) & $w_{k+1}$ \\ \hline
0 & $0$ & $2(w_0-3)-2(\phi_1^{0}-3)+2(\phi_1^{0}-3)=2(w_0-3)=-6$ & $w_1=0-0.1(-6)=0.6$ \\[2mm]
1 & $0$ (still) & $2(w_1-3)-2(\phi_1^{0}-3)+2(\phi_1^{0}-3)=2(w_1-3)=-4.8$ & $w_2=0.6-0.1(-4.8)=1.08$ \\[2mm]
2 & $0$ (still) & $2(w_2-3)=-3.84$ & $w_3=1.08-0.1(-3.84)=1.464$ \\ 
\end{tabular}
\end{center}

Objective values:  

\[
f(w_0)=9,\quad f(w_1)=\bigl(0.6-3\bigr)^2=5.76,\quad 
f(w_2)=\bigl(1.08-3\bigr)^2=3.6864,\quad 
f(w_3)=\bigl(1.464-3\bigr)^2=2.3689.
\]

The iterates move monotonically towards the optimum \(w^{\star}=3\).

\newpage
\section*{Assumptions}
\begin{assumption}[Convexity and Smoothness]\label{ass:convex}
Each component \(f_i\) is convex and \(L\)-smooth, i.e.
\[
f_i(y)\ge f_i(x)+\langle\nabla f_i(x),y-x\rangle,\qquad
\|\nabla f_i(x)-\nabla f_i(y)\|\le L\|x-y\|,
\;\forall x,y\in\R^{d}.
\]
\end{assumption}
\begin{assumption}[Strong Convexity (optional)]\label{ass:strong}
There exists \(\mu>0\) such that for all \(x\),
\[
f(x)-f(x^{\star})\ge\frac{\mu}{2}\|x-x^{\star}\|^{2},
\]
where \(x^{\star}=\arg\min f\).
\end{assumption}
\begin{assumption}[Uniform Sampling]\label{ass:sampling}
At each iteration the index \(i_k\) is drawn independently from the uniform
distribution on \(\{1,\dots ,n\}\).
\end{assumption}
\begin{assumption}[Step Size]\label{ass:stepsize}
The constant step size satisfies  
\[
0<\alpha\le \frac{1}{3L}.
\]
\end{assumption}

\section*{Main Convergence Theorem}
\begin{theorem}[Convergence of SAGA]\label{thm:convergence}
Assume \ref{ass:convex}--\ref{ass:stepsize} hold.
\begin{enumerate}[label=(\alph*)]
\item \textbf{Strongly convex case (\(\mu>0\)).}
Let \(\alpha\le \frac{1}{3L}\). Then for all \(k\ge 0\),

\[
\E\!\bigl[f(x_k)-f(x^{\star})\bigr]
\;\le\;
\bigl(1-\tfrac{\mu\alpha}{2}\bigr)^{k}
\Bigl[f(x_0)-f(x^{\star})+\tfrac{L}{2}\tfrac{1}{n}\sum_{i=1}^{n}\|\phi_i^{0}-x^{\star}\|^{2}\Bigr].
\tag{1}
\]

Consequently \(\E[f(x_k)-f(x^{\star})]=\mathcal{O}\!\bigl((1-\frac{\mu}{6L})^{k}\bigr)\).

\item \textbf{Plain convex case (\(\mu=0\)).}
With the same step size, for any \(K\ge 1\),

\[
\frac{1}{K}\sum_{k=0}^{K-1}\E\!\bigl[f(x_k)-f(x^{\star})\bigr]
\;\le\;
\frac{2L}{K}\Bigl[\|x_0-x^{\star}\|^{2}
      +\tfrac{1}{n}\sum_{i=1}^{n}\|\phi_i^{0}-x^{\star}\|^{2}\Bigr].
\tag{2}
\]

Thus the averaged iterate achieves the optimal \(\mathcal{O}(1/K)\) rate.
\end{enumerate}
\end{theorem}

\section*{Theoretical Properties}
\begin{itemize}
\item \textbf{Variance Reduction.} The control variate
\(\nabla f_{i_k}(\phi_{i_k})-\bar g\) makes the stochastic direction unbiased
and its variance decays as the table entries become closer to the optimum.
\item \textbf{Stability w.r.t.\ \(\alpha\).} Condition \(\alpha\le 1/(3L)\) guarantees that the
Lyapunov function (see Lemma~\ref{lem:lyapunov}) is non‑increasing in expectation.
\item \textbf{Comparison.}
For smooth convex problems SGD with a fixed step size yields
\(\mathcal{O}(1/\sqrt{k})\) suboptimality, while SAGA attains the linear
rate (1) under strong convexity and the optimal \(\mathcal{O}(1/k)\) rate (2) without strong convexity.
\end{itemize}

\section*{Preliminaries (Definitions and Lemmas)}
\begin{definition}[Lyapunov function]
For any iterate \(k\) define  

\[
\Psi_k:=\E\!\Bigl[\|x_k-x^{\star}\|^{2}
      +\frac{1}{n}\sum_{i=1}^{n}\|\phi_i^{k}-x^{\star}\|^{2}\Bigr].
\tag{3}
\]
\end{definition}

\begin{lemma}[Smoothness inequality]\label{lem:smooth}
For any \(L\)-smooth convex \(f_i\),

\[
f_i(y)\le f_i(x)+\langle\nabla f_i(x),y-x\rangle+\frac{L}{2}\|y-x\|^{2},
\qquad\forall x,y.
\tag{4}
\]
\end{lemma}

\begin{lemma}[Unbiasedness of the SAGA direction]\label{lem:unbiased}
Conditioned on the sigma‑algebra \(\mathcal{F}_k:=\sigma\{x_0,\phi_i^{0},\dots ,x_k,\phi_i^{k}\}\),

\[
\E\!\bigl[\,\nabla f_{i_k}(x_k)-\nabla f_{i_k}(\phi_{i_k}^{k})+\bar g^{k}\mid\mathcal{F}_k\bigr]=\nabla f(x_k).
\tag{5}
\]
\end{lemma}
\begin{proof}
Because \(i_k\) is uniform,
\[
\E[\nabla f_{i_k}(x_k)\mid\mathcal{F}_k]=\frac{1}{n}\sum_{i=1}^{n}\nabla f_i(x_k)=\nabla f(x_k).
\]
Similarly
\[
\E[\nabla f_{i_k}(\phi_{i_k}^{k})\mid\mathcal{F}_k]=\frac{1}{n}\sum_{i=1}^{n}\nabla f_i(\phi_i^{k})=\bar g^{k}.
\]
Subtracting gives (5). \(\square\)
\end{proof}

\section*{Main Proof (Step‑by‑Step Derivation)}
We prove Theorem~\ref{thm:convergence}.  
All expectations are taken w.r.t.\ the randomness up to iteration \(k\) unless
otherwise noted.

\paragraph{Step 1 (One‑step expansion).}
Starting from the update rule,
\[
x_{k+1}=x_k-\alpha v_k,\qquad 
v_k:=\nabla f_{i_k}(x_k)-\nabla f_{i_k}(\phi_{i_k}^{k})+\bar g^{k}.
\tag{6}
\]

\[
\|x_{k+1}-x^{\star}\|^{2}
 =\|x_k-x^{\star}\|^{2}
    -2\alpha\langle v_k, x_k-x^{\star}\rangle
    +\alpha^{2}\|v_k\|^{2}.                               \tag{7}
\]

\paragraph{Step 2 (Take conditional expectation).}
Condition on \(\mathcal{F}_k\) and use Lemma~\ref{lem:unbiased}:

\[
\E\!\bigl[\langle v_k, x_k-x^{\star}\rangle\mid\mathcal{F}_k\bigr]
   =\langle \nabla f(x_k),x_k-x^{\star}\rangle.          \tag{8}
\]

\paragraph{Step 3 (Bound the variance term).}
Write \(v_k\) as the sum of an unbiased term plus a zero‑mean correction:

\[
v_k = \nabla f(x_k) + \underbrace{\Bigl(\nabla f_{i_k}(x_k)-\nabla f(x_k)\Bigr)
          -\Bigl(\nabla f_{i_k}(\phi_{i_k}^{k})-\bar g^{k}\Bigr)}_{\delta_k}.
\tag{9}
\]

Hence
\[
\|v_k\|^{2}= \|\nabla f(x_k)\|^{2}+2\langle\nabla f(x_k),\delta_k\rangle+\|\delta_k\|^{2}.
\tag{10}
\]

Taking conditional expectation and noting that \(\E[\delta_k\mid\mathcal{F}_k]=0\),

\[
\E\!\bigl[\|v_k\|^{2}\mid\mathcal{F}_k\bigr]
   =\|\nabla f(x_k)\|^{2}+\E\!\bigl[\|\delta_k\|^{2}\mid\mathcal{F}_k\bigr].
\tag{11}
\]

Now expand \(\delta_k\):

\[
\delta_k=\bigl(\nabla f_{i_k}(x_k)-\nabla f_{i_k}(\phi_{i_k}^{k})\bigr)
        -\bigl(\nabla f_{i_k}(x_k)-\nabla f(x_k)\bigr)
        +\bigl(\bar g^{k}-\nabla f_{i_k}(\phi_{i_k}^{k})\bigr).
\tag{12}
\]

Using the inequality \(\|a+b\|^{2}\le 2\|a\|^{2}+2\|b\|^{2}\) twice gives

\[
\|\delta_k\|^{2}\le 2\underbrace{\|\nabla f_{i_k}(x_k)-\nabla f_{i_k}(\phi_{i_k}^{k})\|^{2}}_{A_k}
               +2\underbrace{\|\nabla f_{i_k}(x_k)-\nabla f(x_k)\|^{2}
                              +\|\bar g^{k}-\nabla f_{i_k}(\phi_{i_k}^{k})\|^{2}}_{B_k}.
\tag{13}
\]

Because sampling is uniform,
\[
\E[A_k\mid\mathcal{F}_k]=\frac{1}{n}\sum_{i=1}^{n}\|\nabla f_i(x_k)-\nabla f_i(\phi_i^{k})\|^{2}.
\tag{14}
\]
Similarly,
\[
\E[B_k\mid\mathcal{F}_k]\le
\frac{1}{n}\sum_{i=1}^{n}\|\nabla f_i(x_k)-\nabla f_i(x^{\star})\|^{2}
+\frac{1}{n}\sum_{i=1}^{n}\|\nabla f_i(\phi_i^{k})-\nabla f_i(x^{\star})\|^{2},
\tag{15}
\]
where we used \(\bar g^{k}=\frac{1}{n}\sum_i\nabla f_i(\phi_i^{k})\) and the convexity of the squared norm.

\paragraph{Step 4 (Apply smoothness).}
From the \(L\)-smoothness of each \(f_i\),

\[
\|\nabla f_i(u)-\nabla f_i(v)\|^{2}\le L^{2}\|u-v\|^{2}.
\tag{16}
\]

Insert (16) into (14)–(15) and combine with (13) to obtain

\[
\E\!\bigl[\|\delta_k\|^{2}\mid\mathcal{F}_k\bigr]
\le 4L^{2}\Bigl(\|x_k-x^{\star}\|^{2}
               +\frac{1}{n}\sum_{i=1}^{n}\|\phi_i^{k}-x^{\star}\|^{2}\Bigr).
\tag{17}
\]

\paragraph{Step 5 (Bounding the one‑step Lyapunov decrease).}
Insert (8), (11) and (17) into (7) and take total expectation:

\[
\begin{aligned}
\E\!\bigl[\|x_{k+1}-x^{\star}\|^{2}\bigr]
&\le \E\!\bigl[\|x_k-x^{\star}\|^{2}\bigr]
   -2\alpha\E\!\bigl[\langle\nabla f(x_k),x_k-x^{\star}\rangle\bigr]   \\
&\qquad +\alpha^{2}\E\!\bigl[\|\nabla f(x_k)\|^{2}\bigr]
   +4\alpha^{2}L^{2}\E\!\Bigl[\|x_k-x^{\star}\|^{2}
          +\frac{1}{n}\sum_{i=1}^{n}\|\phi_i^{k}-x^{\star}\|^{2}\Bigr].
\end{aligned}
\tag{18}
\]

\paragraph{Step 6 (Use convexity of \(f\)).}
For convex \(f\),

\[
\langle\nabla f(x_k),x_k-x^{\star}\rangle\ge f(x_k)-f(x^{\star}).
\tag{19}
\]

Moreover, by \(L\)-smoothness,

\[
\|\nabla f(x_k)\|^{2}\le 2L\bigl(f(x_k)-f(x^{\star})\bigr).
\tag{20}
\]

Insert (19)–(20) into (18):

\[
\begin{aligned}
\E\!\bigl[\|x_{k+1}-x^{\star}\|^{2}\bigr]
&\le \E\!\bigl[\|x_k-x^{\star}\|^{2}\bigr]
   -2\alpha\E\!\bigl[f(x_k)-f(x^{\star})\bigr]   \\
&\quad +2\alpha^{2}L\E\!\bigl[f(x_k)-f(x^{\star})\bigr]
   +4\alpha^{2}L^{2}\E\!\bigl[\|x_k-x^{\star}\|^{2}
          +\tfrac{1}{n}\sum_{i}\|\phi_i^{k}-x^{\star}\|^{2}\bigr].
\end{aligned}
\tag{21}
\]

\paragraph{Step 7 (Memory update contribution).}
When the index \(i_k\) is updated, the table term changes only for that
coordinate:

\[
\|\phi_{i_k}^{k+1}-x^{\star}\|^{2}
 =\|x_k-x^{\star}\|^{2}.
\tag{22}
\]

All other \(\phi_j\) stay unchanged. Taking expectation over the uniform
choice of \(i_k\),

\[
\E\!\Bigl[\tfrac{1}{n}\sum_{i=1}^{n}\|\phi_i^{k+1}-x^{\star}\|^{2}\Bigr]
 =\tfrac{1}{n}\sum_{i=1}^{n}\|\phi_i^{k}-x^{\star}\|^{2}
   +\frac{1}{n}\bigl(\|x_k-x^{\star}\|^{2}
        -\|\phi_{i_k}^{k}-x^{\star}\|^{2}\bigr).
\tag{23}
\]

Since the index is uniform, the last term equals  

\[
\frac{1}{n}\bigl(\|x_k-x^{\star}\|^{2}
        -\tfrac{1}{n}\sum_{i=1}^{n}\|\phi_i^{k}-x^{\star}\|^{2}\bigr).
\tag{24}
\]

Thus

\[
\E\!\Bigl[\tfrac{1}{n}\sum_{i}\|\phi_i^{k+1}-x^{\star}\|^{2}\Bigr]
 =\tfrac{1}{n}\sum_{i}\|\phi_i^{k}-x^{\star}\|^{2}
   +\frac{1}{n}\|x_k-x^{\star}\|^{2}
   -\frac{1}{n^{2}}\sum_{i}\|\phi_i^{k}-x^{\star}\|^{2}.
\tag{25}
\]

\paragraph{Step 8 (Lyapunov recursion).}
Add (21) and (25) and use the definition (3) of \(\Psi_k\). After simplification we obtain

\[
\Psi_{k+1}
\le\Bigl(1+4\alpha^{2}L^{2}+\frac{1}{n}\Bigr)\Psi_k
   -2\alpha\Bigl(1-\alpha L\Bigr)\E\!\bigl[f(x_k)-f(x^{\star})\bigr].
\tag{26}
\]

\paragraph{Step 9 (Choosing \(\alpha\)).}
With the step‑size bound \(\alpha\le \frac{1}{3L}\) we have  

\[
1+4\alpha^{2}L^{2}+\frac{1}{n}\le 1+\frac{4}{9}+\frac{1}{n}
\le 1+\frac{1}{2}\;(\text{for any }n\ge 1),
\]
and more importantly  

\[
2\alpha(1-\alpha L)\ge \alpha.
\tag{27}
\]

Hence (26) simplifies to

\[
\Psi_{k+1}\le\Bigl(1-\frac{\alpha\mu}{2}\Bigr)\Psi_k,
\tag{28}
\]
where in the strongly convex case we used the inequality  

\[
f(x_k)-f(x^{\star})\ge \frac{\mu}{2}\|x_k-x^{\star}\|^{2}
\tag{29}
\]
(see Assumption~\ref{ass:strong}).

\paragraph{Step 10 (Deriving the rate).}
Iterating (28) gives  

\[
\Psi_k\le\Bigl(1-\frac{\alpha\mu}{2}\Bigr)^{k}\Psi_0.
\tag{30}
\]

Since \(\Psi_k\) dominates the suboptimality via smoothness,
\[
f(x_k)-f(x^{\star})\le \frac{L}{2}\|x_k-x^{\star}\|^{2}
\le \frac{L}{2}\Psi_k,
\tag{31}
\]
combining (30) and (31) yields exactly the bound (1) of Theorem~\ref{thm:convergence}.

\paragraph{Step 11 (Convex case).}
When \(\mu=0\), inequality (28) becomes  

\[
\Psi_{k+1}\le\Psi_k-\alpha\,\E\!\bigl[f(x_k)-f(x^{\star})\bigr].
\tag{32}
\]

Summing (32) for \(k=0,\dots ,K-1\) gives  

\[
\alpha\sum_{k=0}^{K-1}\E\!\bigl[f(x_k)-f(x^{\star})\bigr]
\le\Psi_0-\Psi_{K}\le\Psi_0,
\]
or  

\[
\frac{1}{K}\sum_{k=0}^{K-1}\E\!\bigl[f(x_k)-f(x^{\star})\bigr]
\le\frac{\Psi_0}{\alpha K}.
\tag{33}
\]

Using \(\alpha\le 1/(3L)\) and the definition of \(\Psi_0\) yields (2).

\paragraph{Conclusion.}
All steps from (1) to (33) are justified under the assumptions
\ref{ass:convex}–\ref{ass:stepsize}. Hence Theorem~\ref{thm:convergence} holds.
\(\square\)

\section*{Rate Derivation (With Constants)}
\begin{itemize}
\item Strongly convex case: choose \(\alpha=\frac{1}{3L}\).  
Then \(\frac{\alpha\mu}{2}= \frac{\mu}{6L}\) and (1) becomes  

\[
\E[f(x_k)-f(x^{\star})]\le
\Bigl(1-\frac{\mu}{6L}\Bigr)^{k}
\Bigl[f(x_0)-f(x^{\star})+\frac{L}{2n}\sum_{i}\|\phi_i^{0}-x^{\star}\|^{2}\Bigr].
\]

\item Convex case: with the same step size,
\[
\frac{1}{K}\sum_{k=0}^{K-1}\E[f(x_k)-f(x^{\star})]
\le \frac{6L}{K}\Bigl[\|x_0-x^{\star}\|^{2}
   +\frac{1}{n}\sum_{i}\|\phi_i^{0}-x^{\star}\|^{2}\Bigr].
\]
\end{itemize}

\section*{Special Cases}
\begin{enumerate}[label=\arabic*)]
\item \textbf{Exact gradient case (\(n=1\)).}
The table contains the single point \(\phi_1\); the update reduces to
\(x_{k+1}=x_k-\alpha\nabla f(x_k)\), i.e. classic gradient descent,
and the rate collapses to the standard \((1-\mu\alpha)^{k}\) linear rate.
\item \textbf{Large‑scale regime (\(n\gg 1\)).}
The term \(\frac{1}{n}\sum_i\|\phi_i^{k}-x^{\star}\|^{2}\) becomes negligible
as the memory points converge, and the algorithm behaves like a variance‑reduced SGD with essentially the same linear rate.
\item \textbf{Non‑strongly convex but PL‑condition.}
If \(f\) satisfies the Polyak–Łojasiewicz inequality
\( \frac{1}{2}\|\nabla f(x)\|^{2}\ge \mu\bigl(f(x)-f(x^{\star})\bigr) \),
the same algebraic steps lead to a linear convergence bound identical to (1) without requiring convexity.
\end{enumerate}

\end{document}